%%
%% style.tex
%% 
%% Made by Alex Nelson <pqnelson@gmail.com>
%% Login   <alex@lisp>
%% 
%% Started on  2026-04-26T11:05:02-0700
%% Last update 2026-04-26T11:05:02-0700
%% 

\section{Code Style Guide}

\begin{node}
I will probably forget my coding conventions, or develop new
conventions as the need arises. This appendix just reviews the
decisions.

We openly acknowledge a lot of our \SML/ style decisions have been
guided by Isabelle/ML's style guide\footnote{\url{https://isabelle.in.tum.de/library/Doc/Implementation/ML.html}}.
\end{node}

\begin{node}
The \SML/ code may be separated from the literate program, so it is
wise to have coding conventions which are adequate if the code
``stands alone''.
\end{node}

\begin{node}[Guiding principle: portability]
We want to have as many platforms supported as possible. This forces
our hand to support Moscow~ML, which is supported almost
everywhere. However, Moscow ML has peculiar conventions, like placing
signatures in their own ``\texttt{.sig}'' file. We are forced to adopt
these conventions.
\end{node}

\begin{node}[Function comments]
Every function should have a comment with at least its type
signature. Any assumptions about the inputs should be given in a
``\texttt{REQUIRES:}'' preconditions clause in this comment. Any guarantees should
be placed in an ``\texttt{ENSURES:}'' postconditions clause.

Multiply preconditions clauses should be understood as conjoined
together (i.e., they all must be true). It may be more readable to
break out preconditions into multiple clauses (e.g., to stress certain
preconditions). Similarly for postconditions, we may have multiple
clauses and we understand that they are all conjoined together.
\end{node}

\begin{node}[Data structure invariants]
We will also try to document any invariants stipulated about a data
structure somewhere near the module using a comment with a
``\texttt{INVARIANT:}'' clause (or possibly several such clauses,
understood that they are all conjoined together). We will document the
invariant(s) in prose using an invariant enunciation, as we did for
partial maps used in unification~(\S\ref{invariant:unification:cycle-free}).
\end{node}

\begin{node}[Naming conventions]
For the most part,
\begin{enumerate}
\item We organize code around abstract data types
\item Abstract data types are modules (\texttt{structure}s) almost
  always with an abstract ``\texttt{t}'' data type to ``piggieback''
  of the structure name (e.g., \texttt{Env.t}, \texttt{Thm.t},
  \texttt{Term.t}, \texttt{Formula.t}, etc.). The only way to
  construct such types are with the constructors provided.
\item Structures are named using \texttt{PascalCase}---the only
  exception is \texttt{FS0} because $FS_{0}$ uses two capitalized letters
\item Signatures are named using \texttt{PascalCase}
\item Functors are named like structures, but with a prefix
  \texttt{Mk} (e.g., \texttt{MkRedBlackTree}, \texttt{MkAssocList}, etc.).
\item Variables and Functions are named using \texttt{snake\_case}
\item Types are named \emph{either} as \texttt{t} for abstract data
  types, \emph{or} using \texttt{PascalCase}
\item Data constructors are named using \texttt{PascalCase}
\item Exceptions are named using \texttt{PascalCase}
\end{enumerate}
We are essentially growing a collection of ``name parts'', and when we
name things we ``concatenate'' them together (separated by an
underscore [for snake case] or case difference [for Pascal case]).
Studies have indicated people remember ``three or less'' without much
loss of concentration (the ``$7\pm2$'' rule seems to lack empirical grounding),
therefore we try to restrict ourselves to 1--3 ``name parts'' in an
identifier (e.g., \texttt{is\_cyclic} has 2 name parts ``is'' and ``cyclic'',
\texttt{refute} has 1 name part, \texttt{refute\_dnf} has 2 name
parts, \texttt{axiom\_ind\_universe} has 3 name parts).
\end{node}

\begin{node}
Following the ``three or fewer'' rule, we can add up to three
apostrophes to the name of an identifier (e.g., \texttt{fns},
\texttt{fns'}, \texttt{fns''}, \texttt{fns'''}, but not \texttt{fns''''}).
If more is needed, use numeric suffixes (e.g., \texttt{fns},
\texttt{fns1}, \texttt{fns2}, \texttt{fns3}, \texttt{fns4}, etc.).
\end{node}

\begin{node}
Following the Standard ML Basis, we are ``naming similar things similarly''.
A generic ``collection'' abstract data type (like a set) would have
similar functions as the Standard Basis's \texttt{List} module
(\texttt{filter}, \texttt{exists}, \texttt{all}, etc.) with type
signatures resembling the existing ones.
\end{node}

\begin{node}[Specific naming conventions]
There are a few specific naming conventions we follow:
\begin{enumerate}
\item Converting \texttt{Foo.t} types to \texttt{Bar.t} types are done
  with the function \texttt{Foo.to\_bar}
\item Variables which hold a formula are usually abbreviated \texttt{fm}
\item Variables which hold a term are usually abbreviated \texttt{tm}
  (not \texttt{t}, which is the type of an abstract data type).
\item Theorems are named \texttt{thm} or \texttt{th}
\item When dealing with an equation or a binary operator, \texttt{lhs} and \texttt{rhs} may
  be used for identifiers referring to the left-hand side (resp.,
  right-hand side) of the equation (resp., binary operator).
\item Predicates start with ``\texttt{is\_}'' prefix, constructors
  start with ``\texttt{mk\_}'' prefix, and destructors/destructuring
  functions start with ``\texttt{dest\_}'' prefix.
\item Use ``\texttt{eq}'' for testing equality of two objects of the
  same type
\item Use ``\texttt{compare(x,y)}'' for a total ordering where if
  $x<y$ returns \texttt{LESS}, if $x>y$ returns
  \texttt{GREATER}. Require ``\texttt{compare(x,y)}'' to return
  \texttt{EQUAL} if and only if ``\texttt{eq~x~y}'' is true.
\end{enumerate}
\end{node}

\begin{node}[Whitespace]
Indentation is 2-spaces, sometimes 1, rarely 4-spaces. Never use tabs.
Indentation follows a logical format following nesting depth.

Wrap lines at 80 characters long.

For complex expressions consisting of multi-clausal function
definitions, \texttt{handle}, \texttt{case}, and \texttt{let}
expressions are trickier. The multiple clauses of \texttt{fun},
\texttt{fn}, \texttt{handle}, and \texttt{case} get extra indentation
to indicate the nesting clearly.

For multi-clausal \texttt{fun} and \texttt{fn} expressions, the
vertical bar aligns with the ``\texttt{n}'' of \texttt{fun} (resp.,
\texttt{fn}). 

The \texttt{case} expression should be wrapped in a parentheses for
extra clarity. The opening parentheses can be counted as taking one
space of the indentation.
\end{node}

\begin{node}[Indenting modules]
Signatures get indentation after the opening \texttt{sig}.
But neither functors nor structures receive any extra indentation.
\end{node}

\begin{node}
Always end lines with semicolons where sensible. This includes
\texttt{let}-expressions with more than one local declaration, at the
end of a \texttt{val} or \texttt{fun} declaration, at the end of each
declaration in a signature, after the \texttt{end} in a \texttt{let}
or \texttt{local} expression. After each type declaration. After an
exception declaration.

Semicolons end ``sentences'' in \SML/, that is how we use them. Some
people can write 500 page long sentences like James Joyce. But only
Joyce can succeed at doing so.
\end{node}